%
% 3-epsilon.tex
%
% (c) 2026 Prof Dr Andreas Müller
%
\section{Das moderne Fundament}
\kopfrechts{Das moderne Fundament}
Die Entwicklung der Analysis im 17.~und 18.~Jahrhundert fand ohne
sorgfältig gerechtfertigte Grundlagen statt.
Erst im 19.~Jahrhundert war der durch aufgetretene Mehrdeutigkeiten
oder sogar Widersprüche entstandene Leidensdruck gross genug, dass
die Analysis auf ein tragfähiges Fundament solider Definitionen
gestellt wurde.

\subsection{Der Funktionsbegriff
\label{buch:einleitung:fundament:subsection:funktionsbegriff}}
Nicht einmal der Begriff einer stetigen Funktion war präzis definiert.
In seiner \emph{Introductio in analysin infinitorum}
\cite{buch:introductio1,buch:introductio2} beginnt er mit der
Vorstellung, dass Funktionen durch einen ``analytischen'' Ausdruck
aus Variablen und Konstanten gegeben sein.
In dieser Vorstellung ist es auch natürlich, mehrfache Werte 
wie zum Beispiel bei $\pm\!\sqrt{x}$ für $x>0$ zuzulassen.
Später, in seinen \emph{Institutions calculi differentialis}
\cite{buch:institutiones} verwendet er dagegen eine allgemeinere
Definition, die aber eher noch unklarer ist.
Er sagte, eine Funktion sei eine Grösse, die sich ändere, wenn sich die
Variablen, von denen sie abhängt, ändern.

Fourier verstand unter einer Funktion etwas, was sich durch eine
\index{Fourier, Joseph}%
trigonometrische Reihe ausdrücken liess.
\index{trigonometrische Reihe}%
\index{Reihe, trigonometrisch}%
Er postulierte, dass \emph{jede} Funktion durch eine solche Reihe
dargestellt werden kann.
Für viele dieser Funktionen ist es aber nicht mehr möglich, einen
``analytischen'' Ausdruck dafür anzugeben.
Die Entdeckung Fouriers machte erst klar, dass zwischen der von der
Funktion gegebenen Zuordnung von Funktionswerten zu Argumentwerten
und der Darstellung dieser Zuordnung als ``analytischer'' Ausdruck
zu unterscheiden ist.

Auch Cauchy, der viel zur Formalisierung der Analysis beigetragen hat,
\index{Cauchy, Augustin-Louis}%
betrachtete Funktionen im Wesentlichen gleichbedeutend mit
``analytischen'' Ausdrücken. 
Das moderne Konzept einer eindeutigen Zuordnung eines Funktionswertes
$f(x)$ zu jedem Wert $x$ des Definitionsbereichs taucht explizit
erst bei Dirichlet 1837 auf.
Er schreibt sogar explizit, dass es nicht nötig sei für eine Funktion
in diesem Sinne, dass sie durch mathematische Operationen ausgedrückt
werden kann.

%
% Konvergenz und Grenzwert
%
\subsection{Konvergenz und Grenzwert}
Das Konzept eines Grenzwertes hat lange Geschichte.
\index{Grenzwert}%
Schon die in Abschnitt~\ref{buch:einleitung:frueh:subsection:vorlaeufer}
dargestellte Exhaustionsmethode ging von einem intuitiven Verständnis
dafür aus, dass eine Approximation immer genauer werden kann.
\index{Approximation}%
Newton formulierte einen Grenzwert als eine Grösse, zu der der Abstand
\index{Newton, Isaac}%
kleiner als jede gegebene Vorgabe werden kann.

Die moderne Definition der Stetigkeit geht auf Bernard Bolzano
\index{Stetigkeit}%
zurück.
\index{Bolzano, Bernard}%
Sie postuliert, dass ein Funktion im Punkt $x_0$ stetig ist, wenn
es für jedes $\varepsilon>0$ ein $\delta >0$ gibt, so dass
\[
|f(x)-f(x_0)|<\varepsilon
\qquad\text{falls}\qquad
|x-x_0|<\delta.
\]
Sein Werk bleib jedoch kaum beachtet.
Cauchy formalisierte den Begriff des Grenzwertes mit dieser Technik,
\index{Cauchy, Augustin-Louis}%
die heute gemeinhin $(\varepsilon,\delta)$-Definition genannt wird.

\subsection{Integral}
Für stetige Funktionen spielt es nicht so eine grosse Rolle, wie man 
das Integral definiert.
Cauchy hat zum Beispiel das Integrationsintervall in Teilintervalle 
aufgeteilt und daraus Rechtecke gebildet, deren Höhe durch den Funktionswert
am Ende des Teilintervalls gegeben ist.
Für die damals üblichen Standards der Strenge der Beweisführung hat
das ausgereicht.
Später hat sich gezeigt, dass diese Definition für die betrachteten
Funktionen tatsächlich äquivalent zu ausgefeilteren Definition des
Integrals waren.

Die Fourier-Koeffizienten werden mit Hilfe der Fourier-Integrale
berechnet.
Für die Anwendung der Fourier-Theorie auf nicht stetige Funktionen
wurde die Antwort auf die Frage dringen, welche Funktionen ein
wohldefiniertes Integral haben.
Dies war die Motivation für Bernhard Riemann, eine Definition 
\index{Riemann, Bernhard}%
zu formulieren und streng zu beweisen, dass eine ausreichen grosse
Klasse von Funktionen integrierbar im Sinne seiner Definition sind.
Das Riemann-Integral ist sehr intuitiv verständlich und wird daher
\index{Riemann-Integral}%
noch heute im Analysisunterricht eingesetzt.

Das Riemann-Integral hat allerdings eine Reihe von Defiziten.
Zum Beispiel gibt es Funktionen, denen man intuitiv einen Integralwert
zuweisen kann, die aber nach der riemannschen Definition nicht integrierbar
sind.
Dazu gehört die Dirichlet-Funktion
\index{Dirichlet-Funktion}%
\[
D
\colon
\mathbb{R}\to\mathbb{R}
:
x\mapsto\begin{cases}
1&\qquad x\in\mathbb{Q}\\
0&\qquad \text{sonst,}
\end{cases}
\]
welche nur in abzählbar vielen Punkten von 0 verschieden ist und der
man daher den Integralwert $0$ zuteilen möchte.

Besonders für die Anwendung, für die das Riemann-Integral definiert worden
ist, nämlich für die Fourier-Theorie, ist von Bedeutung, dass das
Integral sich mit Grenzwerten von Funktionenfolgen gut verträgt.
Auch hier zeigt sich, dass es konvergente Funktionenfolgen gibt, deren
Integrale nicht konvergieren.

Der Fundamentalsatz der Infinitesimalrechnung ist das Resultat,
dass das Integral ein inverse Operation zur Ableitung einer
Funktion liefert.
Für das riemannsche Integral ist diese Aussage nicht automatisch
war, es sind zusätzliche Anforderungen an den Integranden notwendig.
Man würde sich wünschen, dass alle Integrale stetiger Funktionen
differenzierbar sind und als Ableitung den Integranden ergeben.

Im 19.~Jahrhundert wurden verschiedene Ansätze versucht, um den 
riemannschen Integralbegriff zu verfeinern, die sich jedoch alle als
äquivalent zur ursprünglichen Definition erwiesen und damit auch
deren Defizite teilte.

Erst die Konstruktion von Henri Lebesgue hat diese Problem überwunden.
\index{Lebesgue, Henri}%
Dazu war jedoch die etwas schwerfällige Entwicklung einer Masstheorie
\index{Masstheorie}%
für Teilmengen des Definitionsbereichs der zu integrierenden Funktionen
notwendig.
Dies verbessert die Verträglichkeit mit Grenzwerten und vergrössert
die Klasse der Funktionen, für die der Hauptsatz gilt.
Diese Form der Integrationstheorie ist die Basis mächtiger Erweiterungen,
die das Fundament der modernen Wahrscheinlichkeitstheorie bilden.

Andere moderne Verfeinerungen wie das Henstock-Kurzweil-Integral
finden einen interessanten Mittelweg.
Durch minimale Erweiterungen der riemannschen Konstruktion kann ein
Integral von reellen Funktionen auf $\mathbb{R}$ definiert werden,
\index{R@$\mathbb{R}$}%
welches die guten Eigenschaften des Lebesgue-Integrals auf ein noch
grössere Klasse von Funktionen ausdehnt und auch den Hauptsatz zu
einer Selbstverständlichkeit macht.

%
% Nichtstandard-Analysis
%
\subsection{Nichtstandard-Analysis}
Die Entwicklungen des 19.~Jahrhunders haben die Analysis auf eine
solide Grundlage gestellt, sie aber auch der intuitiven Rechnung
beraubt, die Eulers oft nicht rigorosen Rechnungen so gut
verständlich gemacht haben.
Die Strenge der $\varepsilon$-$\delta$-Technik hat der Analysis die
unendlich kleinen Grössen ausgetrieben, die dafür nötig waren.
Die Nichtstandard-Analysis gibt der Theorie solche Elemente
\index{Nichtstandard-Analysis}%
zurück.

Besonders leicht zugänglich ist die \emph{Internal Set Theory},
\index{internal set theory}%
die ein Modell der reellen Zahlen mit einem zusätzlichen Prädikat
$\operatorname{st}(x)$ ist.
Das Prädikat ist wahr für die vertrauten reellen Zahlen.
In $\mathbb{R}$ gibt es aber auch positive Elemente $h\in\mathbb{R}$, 
die kleiner sind als alle Standardelemente:
\[
0<h<r\qquad\forall r\in\mathbb{R}_+(\operatorname{st}(r)).
\]
Solche Elemente sind aus Sicht der standardreellen Zahlen vernachlässigbar
klein, sie werden auch als unendlich klein oder infinitesimal bezeichnet.
Man kann aber damit wie gewohnt rechnen und erhält damit die Möglichkeit,
die Ableitung wieder rein algebraisch zu finden.
Die Ableitung einer Funktion $f(x)$ ist diejenige reelle Zahl $f'(x)$,
die sich vom Differenzenquotienten mit Differenz $h$ nur um eine 
\index{Differenzenquotient}%
unendlich kleine Zahl unterscheiden:
\[
\biggl|
f'(x)
-
\frac{f(x+h)-f(x)}{h}
\biggr|
<
r
\qquad\forall r\in\mathbb{R}_+(\operatorname{st}(r)).
\]
Damit ist die Intuition der Algebra in die Analysis zurückgekehrt.

Die Nichtstandard-Analysis führt nicht zu neuen Ergebnissen in der
Analysis.
Das sogenannte Transferprinzip besagt nämlich, dass sich für jedes Resultat,
\index{Transferprinzip}%
dass mit der Nichtstandard-Analysis gefunden worden ist, auch ein
klassischer Beweis angeben lässt.

